%% =====================================================================
%%
%% This is file thesis.tex, a template for Tulane University Thesis
%% (PHYSICS DEPT.), Version 1, by Tim Schuler (based on MATH Version 1.2  
%% by Dmitri Alexeev). It is supplied with files TUT.cls (class 
%% definitions) and TUT.bst (bibliography definitions). This is the 
%% only file you need to edit. Simply follow the instructions provided 
%% for each line. TUT.cls replaces the book document class, which means 
%% that you should have no problems with using it if your system works 
%% with book.cls.
%% =====================================================================

% TUT.cls document class automatically
%% =====================================================================
%% TULANE SSE PRINTING GUIDELINES:
%% 1. The final copy must be printed on 8.5 x 11-inch white paper.
%% 2. The paper must be acid-free, 100% cotton fiber, and 20 lb bond.
%% 3. All margins should be 1 inch, EXCEPT chapter titles and the very
%%    first page of the Abstract, TOC, etc., which must be 2 inches from the top.
%%    (This TUT.cls file attempts to handle these margins automatically)
%% =====================================================================

\documentclass{TUT}
\usepackage{graphicx}
\usepackage{placeins}
\usepackage[utf8]{inputenc}
\usepackage{amsmath}
\usepackage{listings}

% --- MODERN PACKAGES INJECTED (2026) ---
\usepackage[hidelinks]{hyperref}   % Clickable links in PDF
\usepackage{float}      % For exact [H] placement of figures/tables
\usepackage{endnotes}   % For endnotes support
\usepackage{glossaries} % For glossary support
\usepackage{imakeidx}   % For index support
% ---------------------------------------

\makeindex
\makenoidxglossaries
\newglossaryentry{math}{
    name=mathematics,
    description={The abstract science of number, quantity, and space}
}
\newtheorem{theorem}{Theorem}

\graphicspath{ {images/} }





% The margins are set correctly by default. However, if your printer
% gives too much of an offset, you may need to define these to
% some proper values depending on your printer.
% Use the file testpage.tex to figure out values for your printer.
%\voffset=-.25in
%\hoffset=0in

% Put your additional packages here.
%\usepackage{pb-diagram,lamsarrow,pb-lams} 
%\usepackage{amstext,amsmath,amssymb,amsfonts,amsthm}
%\dgARROWLENGTH=1.5em

% If you want to have an Index at the end
%\usepackage{makeidx}
%\makeindex

% Doublespacing is required by the Graduate School. You may still
% use singlespacing to view and print your drafts.
\doublespacing
%\singlespacing

% You may want to use these commands (optional)
%\renewcommand{\figurename}{Diagram}
%\renewcommand{\listfigurename}{List of Diagrams}

% If you have a file with your favorite TeX definitions 
% (e.g. mycoms.tex), you may input it here.
%\input{mycoms}

%%%%%%%%%%%%%%%%%%%%%%%%%%%%%%%%%%%%%%%%%%%%%%%%%%%%%%%%%%%%%%%%%%%%%%%%
%%%%%%%%%%%%%%%%%%%%%%%%% Necessary Definitions %%%%%%%%%%%%%%%%%%%%%%%%
%%%%%%%%%%%%%%%%%%%%%%%%%%%%%%%%%%%%%%%%%%%%%%%%%%%%%%%%%%%%%%%%%%%%%%%%

% Your name
\author{Your Name Here}

% Do you work on a thesis or dissertation?
\thesisform{Dissertation}

% The title of your thesis/dissertation
\title{Your Thesis/Dissertation Title Here}

% The date of submission (OBSERVE THE STYLE!)
\degreedate{Date of Submission}

% The year of the submission date
\degreeyear{\the\year}

% If you want the Copyright notice to be printed, give 
% the Copyright year, otherwise comment out or leave the year empty.
\copyrightyear{\the\year}

% Your Department 
\department{Department Of XXX}

% Your School
\school{School of XXX}

% The full (unabbreviated) name of the degree
\degree{Doctor Of Philosophy}

% The name of your committee's chair
\chair{Chair Name Here, Ph.D.}

% The names of your other committee members, one by one.
% You may have up to SIX other members, ordered alphabetically.
\membera{Member A Name, Ph.D.}
\memberb{Member B Name, Ph.D.}
\memberc{Member C Name, Ph.D.}
\memberd{Member D Name, Ph.D.}
\membere{}
\memberf{}

% If you have a member (or yourself) with a very long name, you will
% need to increase the \approvallength. It is set by default to
% 2.5in, which is sufficient in most cases.
%\approvallength=2.5in 


%===============
\begin{document}
%===============


% The abstract environment produces the abstract title page and the
% abstract, which are required for dissertations. If you don't need
% the abstract (and its title page), then comment it out.
%
% IMPORTANT OVERRIDE RULE: 
% The abstract for a Ph.D. dissertation MUST NOT EXCEED 350 WORDS!
% If it exceeds 350 words, University Microfilms will edit it.
\begin{abstract}
  Write your abstract here.
\end{abstract}


% The maketitle command produces the title page for the thesis and
% the copyright notice if the \copyrightyear has been defined above.
\maketitle


%%%%%%%%%%%%%%%%%%%%%%%%%%%%%%%%%%%%%%%%%%%%%%%%%%%%%%%%%%%%%%%%%%%%%%%%%
\chapter*{Acknowledgments}
%%%%%%%%%%%%%%%%%%%%%%%%%%%%%%%%%%%%%%%%%%%%%%%%%%%%%%%%%%%%%%%%%%%%%%%%%
% All pages following the title are supposed to be numbered in roman
% style at the bottom, up to the Introduction.
\pagenumbering{roman}
\pagestyle{plain}
\setcounter{page}{2}

Write your Acknowledgments here.


%%%%%%%%%%%%%%%%%%%%%%%%%%%%%%%%%%%%%%%%%%%%%%%%%%%%%%%%%%%%%%%%%%%%%%%%%
\chapter*{Foreword}
%%%%%%%%%%%%%%%%%%%%%%%%%%%%%%%%%%%%%%%%%%%%%%%%%%%%%%%%%%%%%%%%%%%%%%%%%

Write your Foreword here or comment it out if you don't need it.


%%%%%%%%%%%%%%%%%%%%%%%%%%%%%%%%%%%%%%%%%%%%%%%%%%%%%%%%%%%%%%%%%%%%%%%%%
% Contents
%%%%%%%%%%%%%%%%%%%%%%%%%%%%%%%%%%%%%%%%%%%%%%%%%%%%%%%%%%%%%%%%%%%%%%%%%
% List your contents or figures (and other similar stuff) here.
\tableofcontents
\listoftables
\listoffigures


%%%%%%%%%%%%%%%%%%%%%%%%%%%%%%%%%%%%%%%%%%%%%%%%%%%%%%%%%%%%%%%%%%%%%%%%%
\chapter{Introduction}
%%%%%%%%%%%%%%%%%%%%%%%%%%%%%%%%%%%%%%%%%%%%%%%%%%%%%%%%%%%%%%%%%%%%%%%%%
% Regularly numbered pages start with the Introduction and continue
% throughout the rest of the thesis
\pagestyle{headings}
\pagenumbering{arabic}

% It is recommended that you keep this main file small by inputing
% the chapters using the ``input'' commands. But if your thesis is
% relatively small, you may write it all in this single file.
% The following line assumes that the Introduction is in the file
% intro.tex.

%\par starts a new paragraph and \noindent suppresses default indention. 
\noindent
**********************************************************************\\
This would be your introduction if you did not want to use the input
command.\\
**********************************************************************\\


\input{Intro}

%%%%%%%%%%%%%%%%%%%%%%%%%%%%%%%%%%%%%%%%%%%%%%%%%%%%%%%%%%%%%%%%%%%%%%%%%
\chapter{Chapter on Examples}
%%%%%%%%%%%%%%%%%%%%%%%%%%%%%%%%%%%%%%%%%%%%%%%%%%%%%%%%%%%%%%%%%%%%%%%%%
% Use the same idea with the input commands for all your other chapters. This will load examples.tex.
%% GENERAL GUIDELINE FOR ILLUSTRATIONS, TABLES, AND FIGURES:
%% Illustrations, especially graphs and charts, should be placed as close as possible 
%% to their first references in the text. 
%% If an illustration, table, or figure appears on a text page, three blank lines 
%% should be left above it and three blank lines should be left below it. 
%% (The 'tulane-math-dissertation.sty' file handles this automatically via \intextsep).
RainThis chapter will demonstrate the formatting of tables and figures within the thesis.

\section{Figures}

The Tulane guidelines state that tables and figures should not be embedded inside the text itself, but rather placed as close to their first reference as possible, separated by three blank lines above and below.

\begin{figure}[H]
    \centering
    % In a real dissertation, use \includegraphics{my_figure} here.
    \rule{4cm}{4cm} % A black box standing in for an image
    \caption{An example figure. Notice the spacing above and below this float block.}
    \label{fig:example_fig}
\end{figure}

The spacing for floating environments like Figure \ref{fig:example_fig} is automatically managed by the {\tt intextsep} configuration in our style package.

\section{Tables}

You can similarly insert tables.

\begin{table}[H]
    \centering
    \begin{tabular}{|c|c|c|}
        \hline
        \textbf{Value 1} & \textbf{Value 2} & \textbf{Result} \\
        \hline
        1                & 2                & 3               \\
        4                & 5                & 9               \\
        \hline
    \end{tabular}
    \caption{An example table formatting.}
    \label{tab:example_tab}
\end{table}

The Table \ref{tab:example_tab} has been added automatically to the ``List of Tables" located in the preliminary pages of the dissertation.


Here's another citation example \cite{Chen-Gaffey2015}.


%%%%%%%%%%%%%%%%%%%%%%%%%%%%%%%%%%%%%%%%%%%%%%%%%%%%%%%%%%%%%%%%%%%%%%%%%
\chapter{Conclusion}
%%%%%%%%%%%%%%%%%%%%%%%%%%%%%%%%%%%%%%%%%%%%%%%%%%%%%%%%%%%%%%%%%%%%%%%%%
%\input{valdoms}

Stop when you are done with all the chapters.




%%%%%%%%%%%%%%%%%%%%%%%%%%%%%%%%%%%%%%%%%%%%%%%%%%%%%%%%%%%%%%%%%%%%%%%%%
%%%%%%%%%%%%%%%%%%%%%%%%%%%%%%% References %%%%%%%%%%%%%%%%%%%%%%%%%%%%%%
%%%%%%%%%%%%%%%%%%%%%%%%%%%%%%%%%%%%%%%%%%%%%%%%%%%%%%%%%%%%%%%%%%%%%%%%%
% File TUT.bst contains the bibliography style definitions for BIBTeX.
% File Bib.bib is assumed to contains the bibliography database 
% to be used by BIBTeX. If you use other filename, change the
% bibliography command below accordingly.
% You can change the name of the Bib.bib file, not TUT.bst file.
% To use BIBTeX, UNCOMMENT the following two lines and COMMENT out
% the sample bibliography a few lines below (or delete it). 
% Press F2 + F11 + F2 + F2 to run bibtex
%
% IMPORTANT BIBLIOGRAPHY RULES:
% If you are in the Sciences, your references must be compiled in 
% the exact order they are cited in the text (e.g., using a style like 'unsrt').
% If you are outside the Sciences, they must be formatted alphabetically.
\bibliographystyle{TUT}
\bibliography{bib}

% If you don't want to use BIBTeX, you will have to create the
% bibliography manually. But remember, Graduate School requires the
% items to appear in ORDER OF CITATION, which means you will need to 
% reorder everything if you change your input file(s). Also, 
% there are specific style requirements. Below is a small sample
% which you can use as a template.
% But really, use BIBTeX!!! 
%\begin{thebibliography}{99}
%\bibitem{AndersonFuller92}
%  F. Anderson and K. Fuller,
%  ``Rings and categories of modules'',
%  Springer-Verlag, 1992
%\bibitem{BazzoniFuchsSalce95}
%  S. Bazzoni, L. Fuchs and L. Salce,
%  ``The hierarchy  of uniserial modules over a valuation domain'',
%  \underline{Forum Math.} {\bf 7}, 247-277 (1995)
%\end{thebibliography}


%%%%%%%%%%%%%%%%%%%%%%%%%%%%%%%%%%%%%%%%%%%%%%%%%%%%%%%%%%%%%%%%%%%%%%%%%
% Endnotes, Glossaries, and Index (Optional uncomment to use)
%%%%%%%%%%%%%%%%%%%%%%%%%%%%%%%%%%%%%%%%%%%%%%%%%%%%%%%%%%%%%%%%%%%%%%%%%
%\theendnotes
%\printnoidxglossaries
%\printindex


%%%%%%%%%%%%%%%%%%%%%%%%%%%%%%%%%%%%%%%%%%%%%%%%%%%%%%%%%%%%%%%%%%%%%%%%%
\chapter*{Biography}
%%%%%%%%%%%%%%%%%%%%%%%%%%%%%%%%%%%%%%%%%%%%%%%%%%%%%%%%%%%%%%%%%%%%%%%%%
% Biography needs to fit on one page, so don't get too involved.

Your biography here.

\end{document}

%%%%%%%%%%%%%%%%%%%%%%%%%%%%%%%%%%%%%%%%%%%%%%%%%%%%%%%%%%%%%%%%%%%%%%%%%
%%%%%%%%%%%%%%%%%%%%%%%%%%%%%%%% E N D %%%%%%%%%%%%%%%%%%%%%%%%%%%%%%%%%%
%%%%%%%%%%%%%%%%%%%%%%%%%%%%%%%%%%%%%%%%%%%%%%%%%%%%%%%%%%%%%%%%%%%%%%%%%